\neskyrius{Įvadas}
\gaires{Šis šablonas skirtas produktui kurti: pavyzdžiui, programai, svetainei,
duomenų bazei, įrenginiui ar mokymosi priemonei. Toliau pateikta skyrių sandara yra
siūlomas būdas aprašyti projektą, o ne atskiras NŠA privalomų skyrių sąrašas.
Pakeiskite gaires savo tekstu. Aprašo apimtis be priedų turi būti 2\,000--3\,000 žodžių;
vien puslapių skaičius jos neparodo.}

\noindent\textbf{Problema ir aktualumas.}
\pildyti{Kam ir kokiomis aplinkybėmis kyla problema? Koks dabartinio sprendimo
trūkumas pagrindžia naujo produkto poreikį? Pateikite faktų ir šaltinių.}

\noindent\textbf{Darbo objektas ir produkto paskirtis.}
\pildyti{Koks produktas kuriamas, kas juo naudosis, kokią konkrečią užduotį jis padės atlikti?}

\noindent\textbf{Tikslas.}
\pildyti{Vienu sakiniu įvardykite numatomą rezultatą: sukurti konkretų produktą
konkrečiam poreikiui ir patikrinti jo tinkamumą pagal pasirinktus kriterijus.}

\noindent\textbf{Uždaviniai.}
\gaires{Suformuluokite tiek uždavinių, kiek reikia tikslui pasiekti. Žemiau pateikta
keturių uždavinių schema yra pavyzdys. Kiekvienam vėliau pateikite atskirą išvadą.}
\begin{enumerate}
\item \pildyti{Išanalizuoti poreikį, šaltinius ir analogus, nustatant produkto reikalavimus.}
\item \pildyti{Suprojektuoti ir pagrįsti pasirinkto sprendimo sandarą.}
\item \pildyti{Įgyvendinti numatytas produkto funkcijas.}
\item \pildyti{Patikrinti produktą ir įvertinti, kiek jis atitinka reikalavimus.}
\end{enumerate}

\noindent\textbf{Metodai ir darbo ribos.}
\pildyti{Kokiais būdais analizuosite poreikį, kursite ir tikrinsite produktą?
Ką apima šis darbas ir ko jame nesiekiama? Hipotezę formuluokite tik jei projekte
iš tikrųjų tikrinate prielaidą; produkto kūrime ji nėra automatinė šablono dalis.}

\noindent\textbf{Etika ir savarankiškumas.}
\pildyti{Kaip nurodysite perimtus sprendimus, saugosite dalyvių duomenis ir
atskleisite pagalbą? Grupiniame darbe įvardykite savo atsakomybę.}
